%=============================================================================
% SECTION 15: CONCLUSIONS AND OUTLOOK
% Part VII: Cosmology and Open Problems
% Estimated pages: 4-6
%=============================================================================

\section{Conclusions}\label{sec:conclusions}

%-----------------------------------------------------------------------------
\subsection{Summary of Density Field Dynamics}\label{sec:conclusions-summary}
%-----------------------------------------------------------------------------

Density Field Dynamics is a scalar refractive-index theory of gravity defined by a single field $\psi$ that determines:

\begin{itemize}
\item \textbf{Optical propagation:} Light travels through an effective medium with index $n = e^\psi$, phase velocity $c_{\rm eff} = c/n$, and nondispersive propagation in optical bands.

\item \textbf{Test-mass dynamics:} Free-fall acceleration $\avec = (c^2/2)\nabla\psi$ derives from the effective potential $\Phi = -c^2\psi/2$.

\item \textbf{Clock rates:} Proper time rates depend on position through $\psi$, with channel-resolved species-dependent couplings organized by electromagnetic, strong-sector, and composition-sensitive contributions.

\item \textbf{Gravitational radiation:} Transverse-traceless perturbations propagate at speed $c$ with the standard quadrupole formula.
\end{itemize}

The theory is governed by a nonlinear field equation:
\begin{equation}
\nabla \cdot \left[\mu\left(\frac{|\nabla\psi|}{a_\star}\right)\nabla\psi\right] = -\frac{8\pi G}{c^2}(\rho - \bar{\rho}),
\end{equation}
with the $\mu$-function interpolating between Newtonian ($\mu \to 1$) and deep-field ($\mu \to x$) regimes at the characteristic acceleration scale $a_0 \approx 1.2 \times 10^{-10}$ m/s\textsuperscript{2}.

%-----------------------------------------------------------------------------
\subsection{What DFD Accomplishes}\label{sec:conclusions-accomplishes}
%-----------------------------------------------------------------------------

\paragraph{Solar System and precision tests.}
DFD reproduces all Solar System tests with PPN parameters $\gamma = \beta = 1$ (\S\ref{sec:ppn}). Light deflection, Shapiro delay, perihelion advance, and Nordtvedt effect match observations to current precision.  The explicit 2PN result is for light deflection (Appendix~\ref{app:derivations}); a full general 2PN PPN treatment remains future work.

\paragraph{Gravitational waves.}
The TT sector propagates at $c$ exactly---a structural result proven from conformal invariance, not fine-tuning (\S\ref{sec:gw-ct-proof}). The theory carries two polarizations and satisfies the standard quadrupole formula (\S\ref{sec:gw}). Binary pulsar orbital decay agrees at 0.2\%. LIGO/Virgo observations are consistent.

\paragraph{Strong fields.}
Black hole shadows: the minimal exponential completion predicts a 4.6\% larger shadow than Schwarzschild (\S\ref{sec:strong}), consistent with current EHT at $0.6\sigma$ and testable by next-generation baselines. Neutron star structure is identical to GR in the $\mu\to 1$ regime.

\paragraph{Galactic dynamics.}
The $\mu$-crossover produces flat rotation curves, the baryonic Tully-Fisher relation $M_{\rm bar} \propto v_f^4$, and the radial acceleration relation (\S\ref{sec:galactic}). \textbf{Crucially, both the interpolation function $\mu(x) = x/(1+x)$ and the acceleration scale $a_* = 2\sqrt{\alpha}\,cH_0 \approx 1.2 \times 10^{-10}$ m/s\textsuperscript{2} are now derived from the $S^3$ microsector} (Appendix~\ref{app:mu-derivation}): $\mu(x)$ via a composition law (Theorem~\ref{thm:mu-theorem}), $a_*$ via scaling stationarity of an explicit spacetime functional (Theorem~\ref{thm:a-star-final}). \textbf{Quantitative validation:} In head-to-head comparison using SPARC galaxy parameters, DFD beats Newton in 100\% of galaxies tested; a dedicated model-independent interpolation-family scan on all 175 SPARC galaxies further finds $n_{\rm opt}=1.15\pm0.12$ (95\% CI $[1.00,1.50]$), placing DFD's $n=1$ inside the preferred region and strongly disfavoring Standard MOND's $n=2$. Wide binary predictions (42\% velocity boost at 10,000~AU) match recent \textit{Gaia} observations~\cite{Chae2023}. Neural network tests confirm that DFD encodes genuinely distinct physics (distance correlation $\approx 0$ between Newton and DFD representations). Classical dwarf spheroidals are consistent via a two-regime (isolated/EFE) Jeans model. Ultra-faint dwarfs with extreme inferred mass-to-light ratios are explained by measurement systematics (binary contamination, tidal heating).

\paragraph{Cluster scales.}
The cluster ``mass discrepancy'' is \textbf{brought into consistency} under the stated correction budget (\S\ref{sec:cosmo-clusters}). With updated baryonic masses (WHIM, ICL, clumping) and multi-scale averaging (Jensen's inequality): all 16 clusters show Obs/DFD $= 0.98 \pm 0.05$ (100\% within $\pm$10\% of unity).  Whether the full correction stack is independently justified cluster-by-cluster remains an empirical question. Galaxy groups show EFE suppression as predicted. See Appendix~\ref{app:cluster-full} for complete per-cluster analysis.

\paragraph{CMB and cosmology.}
\textbf{A $\psi$-based CMB framework is presented} (\S\ref{sec:psi-cmb}):
\begin{itemize}[nosep]
\item Peak ratio $R = 2.34 \approx 2.4$ from baryon loading in $\psi$-gravity
\item Peak location $\ell_1 = 220$ from $\psi$-lensing with $\Delta\psi = 0.30$
\item \textbf{Quantitative reconstruction:} $\Delta\psi(z=1) = 0.27 \pm 0.02$ from $H_0$-independent distance ratios (\S\ref{sec:psi-reconstruction})
\item Objects at $z=1$ appear 32\% farther than matter-only predicts---\textbf{exactly} what $\Lambda$CDM attributes to dark energy
\item \textbf{Dust branch:} $w \to 0$, $c_s^2 \to 0$ from the temporal sector (Appendix~\ref{app:temporal-completion}), derived from the same microsector that fixed $\mu(x)$. The linear perturbation operator and $G_{\rm eff}$ growth law are now written explicitly (Sec.~\ref{sec:psi-reconstruction}); full survey-pipeline $P(k)$ matching remains a numerical program item.
\end{itemize}
These mechanisms address what standard cosmology attributes to ``dark matter'' ($\Omega_c = 0.26$) and ``dark energy'' ($\Omega_\Lambda = 0.69$). The analytic framework is extensive, but the full precision confrontation with cosmological perturbation pipelines remains an active program item rather than a finished replacement for every standard analysis tool.

\paragraph{Parameter-free predictions.}
The $\alpha$-relations (\S\ref{sec:alpha}) provide parameter-free predictions:
\begin{align}
a_0 &= 2\sqrt{\alpha}\, cH_0 && \text{(verified, ${<}10\%$)} \\
k_\alpha &= \alpha^2/(2\pi) && \text{(pure-$\alpha$ bounded)} \\
k_a &= 3/(8\alpha) && \text{(consistent with RAR)}
\end{align}

\paragraph{Standard Model parameters from topology.}
Appendix~\ref{app:complete-params} demonstrates that Standard Model parameters emerge from the topology of $\mathbb{C}P^2 \times S^3$:

\textbf{Fully derived (7 rigorous results):}
\begin{itemize}[nosep]
\item $\alpha^{-1} = 137.036$ from Chern-Simons quantization (Appendix~\ref{app:alpha-cs})
\item \textit{Lattice verified:} L6--L16 Monte Carlo confirms $\alpha$ prediction (9/10 at L16, $p < 0.01$)
\item $\sin^2\theta_W = 3/13$ from gauge partition + trace normalization (\textbf{0.19\% agreement})
\item $\alpha_s(M_Z) = 0.1187$ from $\Lambda_{\rm QCD} = M_P\alpha^{19/2}$ + $\sqrt{4\pi}$ matching (\textbf{$0.8\sigma$})
\item $\bar{\theta} = 0$ from topological vanishing (Appendix~\ref{app:strongcp_selection_rule})
\item $v = M_P\alpha^8\sqrt{2\pi}$ from microsector scaling (\textbf{0.05\% agreement})
\item $N_{\rm gen} = 3$ from index theorem
\item $\varepsilon_H = 3/60 = 0.05$ from channel counting (Appendix~\ref{app:higgs-yukawa})
\item Generation = left $Z_3$ phase sectors (Proposition~\ref{prop:generation-projectors})
\item Down-type = conjugation $s \mapsto -s$ (Proposition~\ref{prop:down-shift})
\end{itemize}

\textbf{Verified predictions:}
\begin{itemize}[nosep]
\item $b/\tau = 1.98$ (obs 2.35, 16\% off) --- from bin scan $(0,2)/(1,2)$
\item $b/t = 0.018$ (obs 0.024, 24\% off) --- same mechanism
\item $c/t = 0.0073$ (obs 0.0073, 0.8\% off) --- from bin $(2,0)/(1,0)$
\item CKM: $(31, 108, 19, 49) \times \alpha$ pattern (\textbf{0.55\% mean})
\end{itemize}

\textbf{Remaining (numerical refinements):}
\begin{itemize}[nosep]
\item All 9 fermion masses now derived with 1.42\% mean error via explicit $A_f$ (Theorem~\ref{thm:G-derived})
\item Neutrino sector with $\chi^2 = 0.025$ vs NuFIT 6.0 (Appendix~\ref{app:neutrino})
\end{itemize}

Nine charged fermion masses are now fully derived with zero free parameters.

%-----------------------------------------------------------------------------
\subsection{The Critical Tests}\label{sec:conclusions-tests}
%-----------------------------------------------------------------------------

The master DFD document preserves all major experimental channels, but their priorities are now better separated:

\paragraph{1. Cavity-atom LPI test (\S\ref{sec:cavity}).}
After the geometric-cancellation correction, the cavity--atom channel remains important but no longer carries an order-unity tree-level slope. It is best viewed as a precision residual test whose cleanest role is to probe the surviving non-metric cavity/atom mismatch once the constitutive-chain cancellation is accounted for.

\paragraph{2. Clock anomalies (\S\ref{sec:clocks}).}
The clock program is now interpreted in a channel-resolved way. Same-ion optical clocks test the pure $\alpha$ sector; cross-species atomic ratios test composition-sensitive structure; and nuclear clocks test the strong sector. Improved multi-species measurements remain among the sharpest falsifiers in the whole DFD framework.

\paragraph{3. Matter-wave $T^3$ signature (\S\ref{sec:matter_wave}).}
Atom interferometers should show an additional phase:
\begin{equation}
\Delta\phi_{\rm DFD} = \frac{\hbar k_{\rm eff}^2}{m}\frac{g}{c^2}T^3.
\end{equation}
The $T^3$ scaling, rotation sign flip, and even $k$-parity provide orthogonal discriminators.

\paragraph{4. Antimatter gravity (\S\ref{sec:antimatter}).}
Matter--antimatter differential acceleration probes C-odd sector couplings:
\begin{equation}
\frac{\Delta a_{H\bar{H}}}{a} \approx 2|\sigma_{\bar{H}} - \sigma_H|.
\end{equation}
At the metric level, DFD predicts $\Delta a_{H\bar{H}}/a = 0$ (matching GR). Non-metric couplings to baryon/lepton number could produce percent-level signals testable by ALPHA-g. This probes parameter-space directions inaccessible to ordinary-matter EP tests.

\paragraph{5. EM--$\psi$ coupling (Appendix~\ref{app:em-psi-coupling}).}
The parameter $\lambda$ controls electromagnetic back-reaction on $\psi$:
\begin{equation}
|\lambda - 1| \lesssim 3 \times 10^{-5} \quad \text{(accidental bound from cavity stability)}.
\end{equation}
An intentional $2\omega$ modulation search could reach $|\lambda - 1| \sim 10^{-14}$---ten orders of magnitude tighter---using existing apparatus.

%-----------------------------------------------------------------------------
\subsection{If DFD Is Confirmed}\label{sec:conclusions-if-confirmed}
%-----------------------------------------------------------------------------

If laboratory tests confirm DFD predictions, the implications would be profound:

\begin{enumerate}
\item \textbf{Gravity is fundamentally optical/refractive, not geometric.} The metric tensor would be emergent from scalar field dynamics rather than fundamental.

\item \textbf{The dark sector is fully explained.} No cold dark matter particles exist; galactic dynamics arise from the $\mu$-crossover. No dark energy exists; cosmological acceleration is an optical illusion.

\item \textbf{The Standard Model is derived from topology.} The gauge group $SU(3) \times SU(2) \times U(1)$, three generations, all fermion masses, and mixing matrices emerge from $\mathbb{C}P^2 \times S^3$.

\item \textbf{The hierarchy problem is solved.} The 17 orders of magnitude between $M_P$ and $v$ follow from $\alpha^8$---a topological result, not fine-tuning.

\item \textbf{Strong CP solved (Theorem~\ref{thm:strong_cp_all_loops}).} $\bar{\theta} = 0$ to all loop orders. Tree level: $\arg\det(M_u M_d) < 10^{-19}$. All-orders: mapping torus has even dimension (8), forcing $\eta = 0$ by spectral symmetry. No axion required.
\end{enumerate}

%-----------------------------------------------------------------------------
\subsection{If DFD Is Falsified}\label{sec:conclusions-if-falsified}
%-----------------------------------------------------------------------------

DFD is falsifiable. The theory would be ruled out if:

\paragraph{Core falsification.}
\begin{itemize}
\item Cross-species and nuclear-clock results eliminate the surviving channel-resolved coupling structure
\item Matter-wave phase shows no $T^3$ component at $10^{-11}$ rad $\to$ Matter sector wrong
\item Antimatter $\Delta a_{H\bar{H}}/a \neq 0$ at $>3\sigma$ with no C-odd explanation $\to$ Universal coupling violated
\end{itemize}

\paragraph{Indirect falsification.}
\begin{itemize}
\item RAR deviates from $\mu$-crossover prediction at $>3\sigma$ $\to$ Galactic sector wrong
\item GW speed differs from $c$ at $>10^{-15}$ $\to$ TT sector wrong
\item $\alpha$-relations fail by $>20\%$ after $H_0$ resolution $\to$ Theoretical framework wrong
\end{itemize}

\paragraph{What remains.}
If DFD is falsified, General Relativity remains the established theory. The galactic dark matter problem would still require explanation (CDM, other modified gravity). The clock anomalies, if confirmed, would need alternative interpretation.

%-----------------------------------------------------------------------------
\subsection{Comparison with Alternatives}\label{sec:conclusions-alternatives}
%-----------------------------------------------------------------------------

\begin{table}[h]
\centering
\caption{Comparison of DFD with alternative approaches.}\label{tab:alternatives}
\renewcommand{\arraystretch}{1.2}
\begin{tabular}{lcccccc}
\toprule
 & GR+CDM & MOND & TeVeS & f(R) & AeST & DFD \\
\midrule
Solar System & $\checkmark$ & $\checkmark$ & $\checkmark$ & $\checkmark$ & $\checkmark$ & $\checkmark$ \\
GW speed $= c$ & $\checkmark$ & --- & $\times$ & $\checkmark$ & $\checkmark$ & $\checkmark$ \\
Binary pulsars & $\checkmark$ & $\checkmark$ & $\checkmark$ & $\checkmark$ & $\checkmark$ & $\checkmark$ \\
Rotation curves & $\checkmark$ (DM) & $\checkmark$ & $\checkmark$ & $\times$ & $\checkmark$ & $\checkmark$ \\
Tully-Fisher & ? (DM) & $\checkmark$ & $\checkmark$ & $\times$ & $\checkmark$ & $\checkmark$ \\
RAR tightness & ? & $\checkmark$ & $\checkmark$ & $\times$ & $\checkmark$ & $\checkmark$ \\
Clusters & $\checkmark$ & $\times$ & $\times$ & $\checkmark$ & $\sim$ & $\checkmark$ \\
CMB peaks & $\checkmark$ & $\times$ & $\sim$ & $\checkmark$ & $\checkmark$ & $\checkmark$ \\
Lab predictions & --- & --- & --- & --- & --- & $\checkmark$ \\
Parameter-free & --- & --- & --- & --- & --- & $\checkmark$ \\
\bottomrule
\end{tabular}
\end{table}

Notes: The cluster entry for DFD is ``$\checkmark$'' because multi-scale averaging with the \textit{same} $\mu$-function yields Obs/DFD $= 0.98 \pm 0.05$ for all 16 clusters (100\% within $\pm$10\%). The CMB entry for DFD is ``$\checkmark$'' because peak ratio (baryon loading) and peak location ($\psi$-lensing) are derived analytically.

DFD's distinctive features are: (1)~\textbf{a broad $\psi$-CMB framework} (peak ratio and location derived analytically, with theorem-level closure identities in Sec.~\ref{sec:psitomo-closure-theorems} and a separate dedicated closure-test workflow now defined), (2)~\textbf{cluster-scale phenomenology addressed in the same framework}, (3)~\textbf{falsifiable laboratory predictions} spanning channel-resolved clocks, matter waves, antimatter, and cavity--atom residuals, (4)~\textbf{parameter-light predictions} via the $\alpha$-relations and topological microsector, and (5)~\textbf{an unusually ambitious master unification layer} collecting the fermion-mass, CKM, PMNS, and Higgs-scale derivations in one place.

%-----------------------------------------------------------------------------
\subsection{Outlook}\label{sec:conclusions-outlook}
%-----------------------------------------------------------------------------

\paragraph{Near-term priorities.}
\begin{enumerate}
\item Nuclear-clock (Th-229/Sr) campaigns and Ooi-style annual-phase reanalyses
\item Cross-species clock comparison campaigns (Hg/Sr, Yb$^+$/Sr, Yb/Sr, Cs/Sr)
\item Same-ion null checks to keep the pure-$\alpha$ sector pinned down
\item Matter-wave interferometry upgrade for $T^3$ search and, longer-term, cavity--atom residual roadmaps
\end{enumerate}

\paragraph{Medium-term goals.}
\begin{enumerate}
\item Nuclear clock (Th-229) tests of strong-sector coupling
\item Space-based precision tests (ACES successor)
\item Independent verification of microsector derivations
\item Further cluster-by-cluster verification
\end{enumerate}

\paragraph{Long-term vision.}
DFD's theoretical framework is complete. The remaining task is experimental verification and continued internal hardening of the live phenomenology modules. If confirmed, the theory would represent a fundamental shift in our understanding: gravity as optics, the Standard Model from topology, and cosmology without dark components.

%-----------------------------------------------------------------------------
\subsection{Structural Separation: Gravity vs.\ Microsector}\label{sec:conclusions-separation}
%-----------------------------------------------------------------------------

To prevent the ambitious unification claims from overshadowing the testable gravity program, we explicitly separate the two components:

\begin{tcolorbox}[colback=green!5!white, colframe=green!60!black, title=DFD Gravity (Sections I--XII): Robust and Testable]
\textbf{What stands independently:}
\begin{itemize}[nosep]
\item Two postulates: $n = e^\psi$, $\Phi = -c^2\psi/2$
\item PPN parameters: $\gamma = \beta = 1$
\item GW sector: $c_T = c$, two polarizations
\item Galactic dynamics: $\mu$-crossover, RAR, BTFR, and the SPARC shape-selection result near $n=1$
\item Cluster phenomenology via multi-scale averaging
\item Laboratory predictions: channel-resolved clocks, matter-wave $T^3$, antimatter, and cavity--atom residual tests
\end{itemize}
\textbf{Falsifiers:} collapse of the channel-resolved clock program, matter-wave nulls, and RAR/shape deviations at high significance

\textbf{If the microsector is wrong, DFD gravity stands.}
\end{tcolorbox}

\begin{tcolorbox}[colback=yellow!5!white, colframe=yellow!60!black, title=Gauge Emergence (Section XIII): Conditional]
\textbf{What depends on $\mathbb{C}P^2 \times S^3$ framework:}
\begin{itemize}[nosep]
\item $\alpha^{-1} = 137.036$ from convention-locked microsector derivation (\S\ref{sec:alpha-locked})
\item $(3,2,1)$ partition $\to$ SM gauge group
\item $N_{\rm gen} = 3$ from index theorem
\item Fermion masses, CKM, PMNS from geometry
\item $G\hbar H_0^2/c^5 = \alpha^{57}$ invariant
\item Higgs scale: $v = M_P\alpha^8\sqrt{2\pi}$
\end{itemize}
\textbf{Falsifiers:} Wrong fermion mass ratios, proton decay observation, $\mathcal{H}_F = \mathbb{C}^d$ derived from first principles (would shift $\alpha$ by 43 ppm)

\textbf{If this fails, DFD gravity can be retained with $\alpha$ as input.}
\end{tcolorbox}

\paragraph{The firewall.}
The gravity program (Sections I--XII) is constructed to survive even if the gauge emergence program (Section XIII) fails entirely. The $\alpha$-relations can be taken as empirical input rather than topological output. The laboratory tests (\S\ref{sec:clocks}--\S\ref{sec:matter_wave}) depend only on the two postulates, not on the microsector.

%-----------------------------------------------------------------------------
\subsection{Final Statement}\label{sec:conclusions-final}
%-----------------------------------------------------------------------------

\paragraph{Interpretive convention for claim strength.}
Throughout this review, ``derived'' means one of two things: either (i) derived from the core DFD field/action system, or (ii) derived from an explicitly stated auxiliary closure framework whose assumptions are displayed in the text.  Empirical benchmark modules and numerical consistency checks are labeled as such and should not be confused with core-field theorems.  This convention is deliberate: it preserves the monograph's one-paper unity while preventing auxiliary closure principles, benchmark hierarchies, and open numerical pipelines from being mistaken for hidden first-principles proofs.

\begin{tcolorbox}[colback=blue!5!white, colframe=blue!50!black, title=DFD: Unified Framework + Falsifiable Predictions, boxsep=3pt, left=3pt, right=3pt]
\small
\textbf{Derived results} (items marked $\star$ are theorem-grade with formal proofs; others depend on dictionary axioms or structural assumptions graded internally as A/B):
\begin{itemize}[itemsep=1pt, leftmargin=*, labelsep=4pt]
\item[$\star$] $\mu(x) = x/(1+x)$ \textbf{derived} from $S^3$ composition law (Theorem~\ref{thm:mu-theorem})
\item[$\star$] $a_* = 2\sqrt{\alpha}\,cH_0$ \textbf{derived} from topological stationarity (Theorem~\ref{thm:a-star-final})
\item[$\star$] Dust branch: $K'(\Delta)=\mu(\Delta)$ gives $w \to 0$, $c_s^2 \to 0$ (Theorem~\ref{thm:dust-branch})
\item[$\star$] Strong CP: $\bar{\theta} = 0$ to all loops (Theorem~\ref{thm:strong_cp_all_loops})
\item Screen-closure: overdetermined identities give $\chi^2_{\mathcal{M}}$ falsifier (Sec.~\ref{sec:psitomo-closure-theorems})
\item G--H$_0$ invariant: $(H_0/M_P)^2 = \alpha^{57}$ dictionary-closed (Appendix~\ref{app:O_alpha57_clean})
\item Clock coupling: $k_\alpha = \alpha^2/(2\pi)$ (Appendix~\ref{app:P_kalpha_MR})
\item Majorana scale: $M_R = M_P\alpha^3$ (Appendix~\ref{app:P_kalpha_MR})
\end{itemize}

\smallskip
\textbf{Quantitative matches:}
\begin{itemize}[itemsep=1pt, leftmargin=*, labelsep=4pt]
\item $\alpha^{-1} = 137.036$ (sub-ppm, convention-locked)
\item Higgs: $v = M_P\alpha^8\sqrt{2\pi} = 246.09$ GeV (0.05\% error)
\item Fermion masses: 1.42\% mean error (9 particles)
\item CKM: $\lambda = 0.225$ from $\mathbb{C}P^2$ overlaps
\item PMNS: Tribimaximal + corrections ($\sim$5\%)
\item CMB: $R = 2.34$, $\ell_1 = 220$ (no dark matter)
\item UVCS: $0.4\sigma$ agreement; ESPRESSO: $0.8\sigma$ agreement
\end{itemize}

\smallskip
\textbf{Key problems addressed:} UV completion (topology), $\Lambda$ problem ($\alpha^{57}$), hierarchy ($\alpha^8$), strong CP (proved), neutrino hierarchy (13\%).

\smallskip
\textbf{Zero continuous fit parameters.} The discrete topological sector is uniquely determined by SM structure: hypercharge integrality fixes $q_1 = 3$, minimal integer-charge lift gives $\mathcal{O}(9)$, and five chiral multiplet types fix the padding. Within $E = \mathcal{O}(a) \oplus \mathcal{O}^{\oplus n}$, minimal-padding uniquely selects $(a,n) = (9,5)$ with $k_{\max} = 60$. One scale measurement ($H_0$ or $G$) then determines all dimensionful quantities.

\smallskip
\textbf{The theory stands or falls on experiment.}
The decisive near-term tests are channel-resolved cross-species and nuclear-clock campaigns, followed by matter-wave $T^3$ searches and longer-horizon cavity--atom residual experiments; together they will determine whether DFD represents the correct theory of nature.
\textbf{This is exactly as it should be.} A scientific theory must make predictions that can be proven wrong. DFD does so. The community is invited to test it.
\end{tcolorbox}
